\section{Frequency sensitivity and local excitation requirements}
\label{sec:frequency-information}
The following identities are derived for the transfer functions in this study.
They explain an ideal periodic steady-state experiment. They do not replace
the finite-horizon, initialized simulations or their working covariance models.
For positive frequency $\omega$, let $G=G_n(i\omega)$ and define the complex
log-transfer sensitivity vector $\bm h=G^{-1}\partial G/\partial\bm\eta$,
where $\bm\eta=(\log K,\log\tau,\log\theta)$. Direct differentiation gives
\begin{equation}
 \bm h_n(\omega)=\left(
 1,\;\frac{-i\omega\tau}{1+i\omega\tau},\;
 \frac{-i\omega\theta}{1+i\omega\theta/n}
 \right)^\top.
 \label{eq:frequency-sensitivities}
\end{equation}
For the pure-delay model, the third component is $-i\omega\theta$.
At zero frequency, only the gain sensitivity remains. At low frequency,
\[
 \log G_n(i\omega)=\log K-i\omega(\tau+\theta)
 -\tfrac12\omega^2(\tau^2+\theta^2/n)+O(\omega^3).
\]
Thus the leading phase change depends on the sum of the thermal and delay
lags. Their separation depends on higher-order frequency behavior; a small
low-frequency fit error can coexist with poor precision for each lag.
The $\theta^2/n$ term is absent in the pure-delay expansion.

To separate this local frequency geometry from startup information, consider
a single periodic input $u(t)=u_0+A\cos\omega t$ after transients have decayed.
Let $\bm d=\partial G(i\omega)/\partial\bm\eta$ and suppose the $N$ observed
phases have zero sine/cosine sums, zero sine--cosine cross sum, and equal
sine/cosine sums of squares $N/2$.
Under independent errors of known variance $\sigma^2$, the exact phase-sum
information is
\begin{equation}
 F=\frac{N}{\sigma^2}\left[
 (u_0K)^2\bm e_1\bm e_1^\top+
 \frac{A^2}{2}\operatorname{Re}(\bm d\bm d^{*})\right],
 \label{eq:periodic-information}
\end{equation}
where $*$ denotes conjugate transpose. A single tone without its DC offset
contributes a matrix of rank at most two for the three real parameters.
With nonzero offset and amplitude, the DC term supplies a separate gain direction:
its constant sensitivity is orthogonal to the oscillatory components over these
phases. The determinant that tests independence of the lag
components is
\begin{equation}
 \begin{vmatrix}
 \operatorname{Re}h_\tau&\operatorname{Re}h_\theta\\
 \operatorname{Im}h_\tau&\operatorname{Im}h_\theta
 \end{vmatrix}
 =\frac{\omega^3\tau\theta(\tau-\theta/n)}
 {(1+\omega^2\tau^2)(1+\omega^2\theta^2/n^2)}.
 \label{eq:frequency-lag-area}
\end{equation}
It is nonzero away from $\tau=\theta/n$. In this ideal experiment, a DC
component and one nonzero tone therefore give local full rank at a regular
parameter point. For pure delay the determinant is
$\omega^3\tau^2\theta/(1+\omega^2\tau^2)>0$.
At the finite-chain repeated pole, $h_\theta=nh_\tau$ at every frequency;
adding frequencies cannot remove that first-order rank deficiency.
These are local differential statements. They do not remove the cold one-stage
global pole-label ambiguity, calibrate a confidence interval, or establish an
optimal input under amplitude, sampling and temperature constraints.

The factor $G$ in $\bm d=G\bm h$ also matters: information depends on output
sensitivity, rather than phase sensitivity alone. The time-domain designs
evaluate that output sensitivity with retained observations, startup assumptions
and a specified covariance. Constant-command startup records can carry transient
information that is deliberately absent from the periodic steady-state calculation.
Analytic derivatives, the determinant and the phase-sum information were checked
against parameter perturbations, numerical determinants and explicit real
sensitivity matrices. These checks contain no additional TC-Lab fits.
